\uuid{uwih}
\titre{Zéros de fonction}
\auteur{bodin ; hanani ; mzari}
\organisation{exo7}
\licence{CC BY-NC-SA 4.0}
\datecreate{2020-01-20}
\niveau{L1}
\module{Analyse}
\chapitre{Continuité, limite et étude de fonctions réelles}
\sousChapitre{Continuité : théorie}
\difficulte{3}
\theme{Théorèmes des valeurs intermédiaires, zéros de fonction}
\exo7qcmid{lille-1.continuite.023}

\contenu{
\begin{question}
\qclassification{123.01,123.02}
\qtags{motcle=zéros de fonction}
Soit $f : [a,b] \to \Rr$ une fonction continue (avec $a < b$). Quelles assertions sont vraies ?
\begin{answers}
    \bad{Si $f(a) \cdot f(b) < 0$ et $f$ croissante alors $f$ s'annule une unique fois sur $[a,b]$.}

    \bad{Si $f(a) \cdot f(b) < 0$ et $f$ n'est pas strictement monotone alors $f$ s'annule au moins deux fois sur $[a,b]$.}

    \bad{Si $f(a) \cdot f(b) < 0$ alors $f$ s'annule un nombre fini de fois sur $[a,b]$.}

    \good{Si $f(a) \cdot f(b) < 0$ et $f$ strictement décroissante, alors $f$ s'annule une unique fois sur $[a,b]$.}
\end{answers}
\begin{explanations}
Le théorème des valeurs intermédiaires implique que si $f : [a,b] \to \Rr$ est continue avec $f(a) \cdot f(b) < 0$ alors
il existe une valeur $c \in [a,b]$ telle que $f(c)=0$. Pour avoir l'unicité de ce zéro, il suffit que $f$ soit strictement croissante ou bien strictement décroissante.
\end{explanations}
\end{question}
}
